%%%BLOCK id=009-01 ch=06 sec=功·功的计算 kind=knowledge alt=none cont=none y=0.03-0.30
$W=Fs\cos\theta$\quad $s$为对地位移 % 公式下方有红笔下划线

功\quad $-8\unit{J}>5\unit{J}$

能\quad $-8\unit{J}<5\unit{J}$

定力做功\quad $W=\left\{\begin{array}{l}\bar{F}s=\dfrac{F_0+F_1}{2}\,s\\[4pt] \text{$F$-$s$图面积}\\[4pt] \text{动能定理}\end{array}\right.$ % CHECK: 原稿首字形似“定”（“定力做功”），括号内三种方法属变力求功，疑为“变力做功”，照原样转录

变力做功\quad 动能定理
%%%END
%%%BLOCK id=009-02 ch=06 sec=功率 kind=formula alt=none cont=none y=0.30-0.34
平均功率\quad $\bar{P}=\dfrac{W}{t}$

瞬时功率\quad $P=Fv\cos\alpha$
%%%END
%%%BLOCK id=009-03 ch=06 sec=功·人对自己做功 kind=knowledge alt=none cont=none y=0.33-0.42
人对自己做功（外界不对人做功）

不知去向的能量（转化为内能）

\begin{tabular}{@{}lll@{}}
加速： & $f_{\text{静}}$与$v$同向 & 生物质能$\to$机械能 \\
匀速 & $f_{\text{静}}=0$ & 生物质能\textsubscript{\unclear{（原地起跳）}}$\to$内能 \\
减速 & $f_{\text{静}}$与$v$反向 & 机械能$+$生物质能$\to$内\unclear{能} \\
\end{tabular}
% 原稿此表位于页面右侧；“加速/匀速/减速”三词带下划线；“匀速”行“生物质能”下方有一行小字（形似“（原地起跳）”）；末行箭头后的字被页边截断
%%%END
%%%BLOCK id=009-04 ch=06 sec=功·非质点对象做功 kind=method alt=none cont=none y=0.43-0.56
非质点对象做功，有质量的绳/杆

\hspace{4em}分段$+$找重心，$W=m_1g\Delta h_1+m_2g\Delta h_2$

%FIG p009_a 0.065 0.498 0.19 0.553
\notefig[0.25]{figures/p009_a.png}
%%%END
%%%BLOCK id=009-05 ch=06 sec=弹簧模型·能量守恒 kind=method alt=07 cont=none y=0.56-0.67
过程写动能\qquad 过程写动量

弹簧\quad $\left\{\begin{array}{l}\pm W_{\text{弹}}=E_{p\text{初}}-E_{p\text{末}}\Rightarrow\text{过程写动能/能量}\\[4pt] F=k\Delta x\to\text{状态写牛二}\end{array}\right.$

能量守恒\quad $\Delta E_{\text{增}}=\Delta E_{\text{减}}$（定量）\quad $E_1+\cdots E_n=$定值（定性）
%%%END
%%%BLOCK id=009-06 ch=06 sec=机车启动 kind=method alt=03 cont=none y=0.69-1.00
状态写牛二\quad $F_{\text{牵}}-f=ma$

牵引力/功率\quad $P=F_{\text{牵}}v=fv_{\text{匀}}$

匀速速率\quad $V_{\text{均}}=\dfrac{P_{\text{额}}}{f}$ % CHECK: 速度下标手写形似“均”，按物理应为“匀”，照原样转录

%FIG p009_b 0.012 0.803 0.865 0.999
\notefig[0.85]{figures/p009_b.png}
% 图 p009_b 包含：左半“恒P”（P-t、v-t、F牵-1/v、a-1/v四图）与右半“恒a”（P-t、v-t、F牵-1/v、a-1/v四图）及其中的标注；图下方两条独立公式如下

$\dfrac{P_0}{v}-f=ma$\qquad $v_1=\dfrac{P_0}{ma+f}$
%%%END
%%%BLOCK id=010-01 ch=06 sec=机车启动·动能定理 kind=formula alt=none cont=none y=0.02-0.11
$t$\ $S$\ 互求\quad 动能定理

\[
Pt-fS=\Delta E_k=\frac{1}{2}mV_m^2-0
\]
\[
P_0t-fS=\frac{1}{2}mV_m^2-\frac{1}{2}mV_1^2
\]
%%%END
%%%BLOCK id=010-02 ch=06 sec=机车启动·突变问题 kind=method alt=03 cont=none y=0.12-0.25
\textbf{突变问题}

加速\quad $\overset{\downarrow}{a}=\dfrac{F-f}{m}=\dfrac{\frac{P}{v}-f}{m}$\quad $\left\{\begin{array}{l}\text{功率突变}\\ \text{阻力突变}\end{array}\right.$\ $v$\ $\left\{\begin{array}{l}\text{加速}\\ \text{减速}\end{array}\right.$

减速\quad $\overset{\downarrow}{a}=\dfrac{f-\frac{P}{v}}{m}$
% “a”上方有一个小符号（形似“↓”或“v”）；“减速”行分子第二项手写形似“P/v”；两行右侧各有一条短弧线
%%%END
%%%BLOCK id=010-05 ch=06 sec=圆弧摩擦·凹凸轨 kind=knowledge alt=04 cont=none y=0.68-0.80
圆弧摩擦：槽大\ 拱小\quad 摩擦多

凹凸轨：\ $t_{\text{槽}}<t_{\text{平}}<t_{\text{拱}}$

%FIG p010_d 0.05 0.7445 0.185 0.793
\notefig[0.25]{figures/p010_d.png}
%%%END
%%%BLOCK id=040-07 ch=06 sec=连续体功率·风力发电机 kind=derivation alt=07 cont=none y=0.86-1.0
\textbf{风力发电机}
%FIG p040_d 0.162 0.89 0.25 0.96
\notefig[0.25]{figures/p040_d.png}
\[
P=\frac{\frac12\rho\,\Delta l\,\pi r^2V^2}{\Delta t}=\frac12\rho\pi r^2V^3 % 原稿分子中的 Δl 与分母 Δt 各被圈出(Δl/Δt=V)
\]
%%%END
%%%BLOCK id=041-09 ch=06 sec=功·合力做功 kind=derivation alt=02 cont=none y=0.83-0.93
%FIG p041_i 0.07 0.84 0.24 0.915
\notefig[0.25]{figures/p041_i.png}
% 图内标注：左侧方块，经绳连到小圆(滑轮/动点)，力 $F$ 斜向上，两侧夹角各标 $\theta/2$，虚线延长，竖直线为边界
\[
F_{\text{合}}=2F\cos\frac{\theta}{2}\qquad
W=2F\cos\frac{\theta}{2}\cdot S\cdot\cos\frac{\theta}{2}=\red{\underline{2FS\cos^{2}\frac{\theta}{2}}}=\red{\underline{FS(1+\sin\theta)}}
\]% 后两项下方各有红色下划线；CHECK: 2cos^2(θ/2)=1+cosθ，原稿末项写作 1+sinθ(笔误或误读)，照原样转录；“cos”后上标写作“2θ”分母“2”，按 cos^2(θ/2) 转录
%%%END
%%%BLOCK id=062-01 ch=06 sec=功和功率 kind=formula alt=none cont=none y=0.05-0.22
\textbf{机械能}
\begin{itemize}
\item 功\quad $W=Fs\cos\alpha$
  \begin{itemize}
  \item $0\le\alpha<90^\circ$\quad 正功
  \item $90^\circ<\alpha\le180^\circ$\quad 负功
  \item $\alpha=90^\circ$\quad 不做功
  \item $W=W_1+W_2+\cdots+\cdots$ % 原稿末尾为省略号
  \end{itemize}
\item 功率\quad $P=Fv\cos\alpha$\qquad $\bar{P}=\dfrac{W}{t}$\qquad $P_{\text{瞬}}=Fv_{\text{瞬}}\cos\alpha$.
\end{itemize}
%%%END
%%%BLOCK id=062-02 ch=06 sec=功和功率 kind=summary alt=none cont=none y=0.23-0.47
\textbf{功和功率}
\begin{itemize}
\item 功的正负判断和恒力做功的计算.
  \begin{itemize}
  \item 力与位移夹角、力与速度夹角
  \end{itemize}
\item 变力做功的五种计算方法.
  \begin{itemize}
  \item 微元法(拉磨)\quad 平均力法(线性变力 $kx$)\quad 等效代换法\quad $F$-$x$图像\unclear{法}\quad 动能定理求变力做功 % CHECK: “图像法”的末字形似“波”，按语义读作“法”
  \end{itemize}
\item 功率的分析与计算
\item 机车启动问题.\quad 方程.
\end{itemize}
%%%END
%%%BLOCK id=062-03 ch=06 sec=动能定理 kind=summary alt=none cont=none y=0.50-0.80
\textbf{动能定理及应用}
\begin{itemize}
\item 动能\quad $E_k=\dfrac{1}{2}mv^2$\quad 只有正值
\item 动能定理\quad $W_{\text{合}}=\Delta E_k=E_{k2}-E_{k1}$
\item 动能定理的理解与简单应用
\item 动能定理结合图像问题
  \begin{itemize}
  \item $F$-$x$图 % CHECK: 原稿 F 与 x 之间的横线似为双横线（“F=x图”），按 F-x 图转录
  \item $E_k$-$x$图斜率为合外力，\unclear{$E_{\text{机}}$}$\text{-}x$斜率为其他力 % 原稿“斜率”二字为小字，写在“为”字上方（插入）；CHECK: “E机”首字形似 F，按物理意义（机械能-位移图斜率为其他力）读作 E
  \end{itemize}
\item 动能定理解决多过程问题\quad 只管功和两点速度
\end{itemize}
%%%END
%%%BLOCK id=062-04 ch=06 sec=机械能守恒 kind=summary alt=none cont=none y=0.82-1.00
\textbf{机械能的守恒及应用}
\begin{itemize}
\item 重力做功与重力势能\quad 保守力
\item 弹性势能.
\item 机械能守恒的判断\quad $W_{\text{其他}}=0$.
\item 单个物体的机械能守恒\quad $\Delta E_K=-\Delta E_p$
\item \unclear{多}个物体的机械能守恒\quad (单体不守恒) % CHECK: 照片底边裁切，最后一行只露出上半，字迹据残笔辨认
\end{itemize}
%%%END
%%%BLOCK id=063-01 ch=06 sec=功能关系 kind=summary alt=none cont=none y=0.10-0.22
\textbf{功能关系\quad 能量守恒定律}
\begin{itemize}
\item 功能关系
\item 能量守恒定律.
\end{itemize}
%%%END
%%%BLOCK id=063-02 ch=06 sec=功能关系 kind=formula alt=10 cont=none y=0.23-0.37
\begin{itemize}
\item 功能关系的理解与应用
  \begin{align*}
  W_{G}&=-\Delta E_p \\ % CHECK: 下标字形似 a/G，按重力功 G 转录
  W_{\text{弹}}&=-\Delta E_p \\
  W_{\text{合}}&=\Delta E_k \\
  W_{\text{其他}}&=\Delta E_{\text{机}} \\
  W_{\text{电}}&=-\Delta E_p
  \end{align*}
\item 摩擦力做功与能量变化的关系\quad 耗散力
  \begin{itemize}
  \item $Q=fs_{\text{相对路程}}$.\quad $W_{\text{克安}}=W_{\text{电}}$.
  \end{itemize}
\end{itemize}
%%%END
%%%BLOCK id=063-03 ch=06 sec=能量守恒定律 kind=knowledge alt=none cont=none y=0.37-0.43
\textbf{能量守恒定律}
\begin{itemize}
\item $\Delta E_{\text{减}}=\Delta E_{\text{增}}$\quad 分段、分模型、分过程
\end{itemize}
%%%END
%%%BLOCK id=075-01 ch=06 sec=功能极值·杆连接体 kind=problem alt=none cont=none y=0.06-0.27
% 页面上方露出的封面及右上角相邻页的倾斜小字（「ρ=1.429 g·L^-1」「…5.72」）属杂物/相邻页，未转录
\textbf{功能极值}

\begin{problem}[1、]
%FIG p075_a 0.03 0.08 0.33 0.18
\notefig[0.4]{figures/p075_a.png}
\begin{solution}
\[
2mg\cdot 2L\sin\theta-3mgL(1-\cos\theta)=\frac{1}{2}\cdot 2m(2V)^{2}+\frac{1}{2}\cdot 3mV^{2}
\]
\[
4\sin\theta+3\cos\theta=5\sin(\theta+37^\circ)\qquad \text{当 }\theta=53^\circ\text{ 时}
\]
此时 $V_B=\sqrt{\dfrac{4}{11}gL}$

\medskip
$V=\omega L$ % CHECK: 手写像「V=WL」，按 ω 转录
\\
$h_B=L(1-\cos\theta)$.
\end{solution}
\end{problem}
%%%END
%%%BLOCK id=075-02 ch=06 sec=功能极值·杆连接体 kind=problem alt=04 cont=none y=0.27-0.52
\begin{problem}[2、]
%FIG p075_b 0.08 0.275 0.375 0.42
\notefig[0.35]{figures/p075_b.png}
起始 $AB$ 与地成 $60^\circ$ 角，$\theta$ 为何值时 $A$ 脱离墙.

$B$ 速度最大时 $A$ 脱离墙
\begin{solution}
\[
mgL(\sin60^\circ-\sin\theta)=\frac{1}{2}m(V_A^{2}+V_B^{2})
\]
\[
V_A\sin\theta=V_B\cos\theta
\]
\begin{align*}
V_B^{2}&=2gL(\sin60^\circ-\sin\theta)\,\frac{1}{1+\cot^{2}\theta}\\
&=2gL(\sin60^\circ-\sin\theta)\sin^{2}\theta\\
&=gL(2\sin60^\circ-2\sin\theta)\sin\theta\,\sin\theta
\end{align*}
当 $\sin\theta=\dfrac{\sqrt{3}}{3}$ 时 $V_B=\dfrac{\sqrt{\sqrt{3}\,gL}}{3}$ 最大.
\end{solution}
\end{problem}
%%%END
%%%BLOCK id=075-03 ch=06 sec=功能极值·杆连接体 kind=problem alt=04 cont=none y=0.52-0.82
\begin{problem}[3、]
%FIG p075_c 0.06 0.53 0.19 0.665
\notefig[0.2]{figures/p075_c.png}
$A$ 在管内 $B$ 在管外\quad 扰动后.
\begin{itemize}
\item[(1)] 求$A$落地速度
\item[(2)] $A$机械能最小时，地对$B$支持力
\item[(3)] 若$m_A=m_B$，$A$机械能最小时 杆与竖直方向夹角的余弦值为
\end{itemize}
\begin{solution}
\textbf{(1)}\quad $m_AgL=\dfrac{1}{2}m_AV_A^{2}$\quad $V_A=\sqrt{2gL}$\quad $V_B=0$

\textbf{(2)}\quad $E_{A\text{机}}+E_{B\text{机}}=m_AgL$

$F_{\text{杆}}=0$ 时 $E_{B\text{机}}$ 最大

$F_{\text{地}}=m_Bg$

%FIG p075_d 0.51 0.71 0.68 0.835
\notefig[0.18]{figures/p075_d.png}

\textbf{(3)}
\begin{align*}
&mgL(1-\cos\theta)=\frac{1}{2}m(V_A^{2}+V_B^{2})\\
&V_A\cos\theta=V_B\sin\theta\\
&V_B^{2}=2gL(1-\cos\theta)\cos^{2}\theta\\
&\phantom{V_B^{2}}=gL(2-2\cos\theta)\cos^{2}\theta
\end{align*}
$\cos\theta=\dfrac{2}{3}$ 时 $V_B^{2}=\dfrac{8}{27}gL$ 最大

$E_{A\text{机}}=mgL-\dfrac{1}{2}mV_B^{2}=\dfrac{23}{27}mgL$
\end{solution}
\end{problem}
%%%END
%%%BLOCK id=075-04 ch=06 sec=功能极值·杆连接体 kind=problem alt=04 cont=none y=0.80-1.00
\begin{problem}[4、]
%FIG p075_e 0.07 0.79 0.26 0.90
\notefig[0.25]{figures/p075_e.png}
\begin{solution}
杆对$b$先做正功后做负功.

对$a$先做负功再做正功.

$a$落地时 $V=\sqrt{2gh}$

$E_{B\text{机}}\,\max=\dfrac{4}{27}mgh$\quad 此时，$a$的加速度为$g$，$F_{\text{杆}}=0$.

$E_{A\text{机}}\,\min=\dfrac{23}{27}mgh$
\end{solution}
\end{problem}
%%%END
%%%BLOCK id=077-01 ch=06 sec=动能定理与机械能守恒·系统 kind=formula alt=none cont=none y=0.07-0.20
\textbf{系统动能定理．}
\[
\sum W_{\text{外}}+\sum W_{\text{内}}=\Delta E_k
\qquad
\text{内力}
\begin{cases}
W\ \text{绳、杆、}f_{\text{静}}\ \text{抵消}\\
W\ \red{\underline{\text{弹性绳}}}\text{、}\red{\underline{\text{弹簧}}}\text{、}\red{\underline{\text{不抵消}}}
\end{cases}
\]
% 原稿：第二行“弹性绳”“弹簧”“不抵消”下为红笔下划线（字为黑色）
\textbf{机械能守恒}
\[
W_{\text{其}}+W_G+W_{\text{弹}}=\Delta E_k
\qquad
W_{\text{其}}=\Delta E_{k}+\Delta E_{pG}+\Delta E_{p\text{弹}}=\Delta E_{\text{机}}
\]
% CHECK: 第二式 ΔE_k 后手写多一小笔（像逗号或下标 1）
\[
W_{\text{其}}=0\ \text{时 机械能守恒}
\]
%%%END
%%%BLOCK id=077-02 ch=06 sec=连接体·能量 kind=method alt=03 cont=none y=0.19-0.31
\textbf{连接体}
\[
\begin{cases}
\text{对1 由动能定理}\\
\text{对2 由动能定理}\\
\text{速度关联}
\end{cases}
\qquad
\fcolorbox{red!75!black}{white}{\begin{tabular}{@{}c@{}}单体不守恒\\ 系统\,才守恒\end{tabular}}
\]
% 上面“单体不守恒/系统才守恒”外为红圈；下面三行外有红笔竖括号
\begin{itemize}
\item 过程为变加变减速，用超失重分析（对××力与 $mg$ 比较）
\item $V_1\ \overset{\text{\scriptsize 正功}}{0}\ \text{有}\ \overset{\text{\scriptsize 负功}}{0}$
\item $V_2\ \text{有}\ \overset{\text{\scriptsize 负功}}{0}\ \overset{\text{\scriptsize 正功}}{\text{有}}$
\end{itemize}
% 小字“正功/负功”写在对应符号上方；第一行正功在第一个0上、负功在最后的0上；第二行负功在“有”“0”之间偏0、正功在最后的“有”上
弹簧速度不能分解
%%%END
%%%BLOCK id=077-03 ch=06 sec=解题分类·过程与方程 kind=summary alt=none cont=none y=0.31-0.42
不以模型分过程而以方程分过程．
\[
\begin{cases}
\text{全过程类}\\
\text{上抛下落类}\\
\text{绳杆连接体类．}
\end{cases}
\]
%%%END
%%%BLOCK id=077-04 ch=06 sec=上抛下落类·双向性 kind=method alt=none cont=none y=0.45-0.62
\textbf{上抛下落类}\quad （有双向性）．
\begin{keybox}
结论\quad 无论 $f$ 是否存在
\end{keybox}
\red{能量 $\times n$\quad 速度 $\times\sqrt{n}$}
% 以下两行为原稿右侧的黑字小注
\[
E_{k0}\times n\quad h\times n\quad E_{k\text{返}}\times n\quad W_f\times n\quad W_G\times n
\]
\[
V_{n\text{返}}\times\sqrt{n}
\]
% CHECK: “返”字辨认存疑（E_k返、V_n返）
\begin{align*}
\text{上}\quad & 0-\tfrac12 mV_0^2=-mgh-fh \quad\text{①}\\
\text{下}\quad & \tfrac12 mV^2-0=mgh-fh \quad\text{②}\\
& \frac{\text{①}}{\text{②}}=\frac{V_0^2}{V^2}=\frac{f+mg}{mg-f}\\
& \tfrac12 mv^2-\tfrac12 mv_0^2=-2hf
\end{align*}
%%%END
%%%BLOCK id=077-05 ch=06 sec=上抛下落类·双向性 kind=method alt=02 cont=none y=0.62-0.80
\red{$0-V_{\text{上}}^2=-g\sin\theta-\mu g\cos\theta$}\\
\red{$V_{\text{下}}^2=g\sin\theta-\mu g\cos\theta$}
% CHECK: 红字 V 的下标手写像 E/F，按上下文取 上/下
%FIG p077_a 0.065 0.72 0.205 0.78
\notefig[0.25]{figures/p077_a.png}
\[
\begin{cases}
W_G = mg\sin\theta\,S\\
W_f = \mu mg\cos\theta\,S
\end{cases}
\qquad
\frac{W_G}{W_f}=\frac{\tan\theta}{\mu}\ \text{为定值}
\qquad
W_f\propto W_G
\]
%%%END
%%%BLOCK id=077-06 ch=06 sec=上抛下落类·双向性 kind=problem alt=none cont=none y=0.81-0.99
\begin{problem}
\unclear{小}球竖直上抛，最高 $H$，$f$ 恒定，上升至 $h$ 处，$E_k=2E_p$，下落至离地 $h$ 处，$E_p=2E_k$，求 $h$。
\begin{solution}
%FIG p077_b 0.0 0.86 0.28 0.97
\notefig[0.3]{figures/p077_b.png}
\[
\left\{
\begin{array}{l}
\uparrow\ E_k=2mgh\qquad \downarrow\ E_k'=\tfrac12 mgh\\[6pt]
h\ \left\{
\begin{array}{l}
-\tfrac32 mgh=-W_f\times 2\\
W_f=-\tfrac34 mgh\\
0-2mgh=-\tfrac34 mgh-mg(H-h)\\
\therefore h=\tfrac49 H
\end{array}
\right.
\end{array}
\right.
\]
% 原稿左侧示意：顶部一条横线，↑自小圆圈向上、↓向下止于小圆圈（上升/下落经过 h 处）
% CHECK: −(3/2)mgh=−W_f×2 与 W_f=−(3/4)mgh 符号不一致（原稿如此）
\end{solution}
\end{problem}
%%%END
%%%BLOCK id=078-01 ch=06 sec=连接体·能量 kind=problem alt=none cont=none y=0.06-0.20
\textbf{连接体}
\begin{problem}[①]
求 $M$ 到碗底速度
%FIG p078_a 0.08 0.083 0.24 0.165
\notefig[0.28]{figures/p078_a.png}
\begin{solution}
\[
\begin{cases}
MgR-W_T=\tfrac12 MV_M^2\\
W_T-\sqrt2\,mgR=\tfrac12 mV_m^2\\
\therefore V_M\cos45^\circ=V_m
\end{cases}
\]
% CHECK: 速度下标 M/m 手写难分，按物理含义取 V_M（M 的速度）与 V_m
\end{solution}
\end{problem}
%%%END
%%%BLOCK id=078-02 ch=06 sec=连接体·能量 kind=problem alt=04 cont=none y=0.06-0.22
\begin{problem}[②]
求 $M$ 到顶端时压力
%FIG p078_b 0.55 0.074 0.71 0.205
\notefig[0.3]{figures/p078_b.png}
\begin{solution}
\[
\begin{cases}
Mg\,\dfrac{2\pi R}{4}-W_T=\tfrac12 MV_M^2\\[4pt]
W_T-mgR=\tfrac12 mV_m^2\\[2pt]
V_M=V_m\\
mg-N=\dfrac{mV_m^2}{R}
\end{cases}
\]
% CHECK: 第1式分子手写像 2πR，其后 −W_T 手写像 −N_T；M/m 下标手写难分
\end{solution}
\end{problem}
%%%END
%%%BLOCK id=078-03 ch=06 sec=连接体·能量 kind=problem alt=none cont=none y=0.21-0.37
\begin{problem}
%FIG p078_c 0.06 0.225 0.335 0.29
\notefig[0.3]{figures/p078_c.png}
$B$ 到地上\\
$A$ 上 $B$ 下，$A$ 刚好能到顶点\\
原 $B$ 在顶点到现在\\
\unclear{提}住 $B$ 不动 $A$ 恰好到顶点
\begin{solution}
\[
\begin{cases}
A:\ W_T-m_1g\sin\theta\,\dfrac{H}{2}=\tfrac12 m_1V_1^2\\[4pt]
B:\ m_2g\,\dfrac{H}{2}-W_T=\tfrac12 m_2V_2^2\\[4pt]
V_1=V_2
\end{cases}
\]
\[
-m_1g\sin\theta\,\dfrac{H}{2}=0-\tfrac12 m_1V_1^2
\]
\end{solution}
\end{problem}
%%%END
%%%BLOCK id=078-04 ch=06 sec=连接体·能量 kind=problem alt=none cont=none y=0.39-0.53
\begin{problem}
\underline{只挂 $C$}\ $B$ 刚离地，把 $C$ 换为 $D$ 求 $B$ 离地速度
%FIG p078_d 0.05 0.41 0.245 0.52
\notefig[0.25]{figures/p078_d.png}
\begin{solution}
\[
\underline{\text{过程}}\text{一}\quad
\begin{cases}
C:\ m_3gh-W_T=0\\
A:\ W_T-m_1gh-E_{p\text{弹}}=0
\end{cases}
\]
\[
\underline{\text{过程}}\text{二}\quad
\begin{cases}
D:\ (m_1+m_3)gh-W_T'=\tfrac12(m_1+m_3)V^2\\
A:\ W_T'-m_1gh-E_{p2\text{弹}}=\tfrac12 m_1V^2
\end{cases}
\]
% CHECK: 过程二中 E_p 与“弹”之间有一字符，像 2 或 3
\textbf{差值法}\quad $m_1gh=\tfrac12(m_1+m_1+m_3)V^2$
% 原稿位置：右上，与“过程一”同高
\end{solution}
\end{problem}
%%%END
%%%BLOCK id=078-05 ch=06 sec=连接体·能量 kind=method alt=none cont=none y=0.56-0.64
%FIG p078_e 0.085 0.578 0.19 0.635
\notefig[0.25]{figures/p078_e.png}
\[
\begin{cases}
W_{G\text{物}}-W_T=\Delta E_{k\text{物}}\\
W_T+W_f+W_{G\text{绳}}=\Delta E_{k\text{绳}}.
\end{cases}
\]
%%%END
%%%BLOCK id=078-06 ch=06 sec=功能关系·系统与单体机械能 kind=problem alt=none cont=none y=0.66-0.79
\begin{problem}
%FIG p078_f 0.07 0.672 0.205 0.732
\notefig[0.28]{figures/p078_f.png}
升降机由静止开始加速上升 $h$．
\begin{solution}
\begin{align*}
\Delta E_{\text{机}B} &= W_N+W_{T\,\text{弹簧对}B}\\
\Delta E_{\text{机系}} &= W_N+W_{T\,A\text{对弹簧}}
\end{align*}
% CHECK: W_T 下标开头有一笔像 3
\end{solution}
\end{problem}
%%%END
%%%BLOCK id=078-07 ch=06 sec=功能关系·系统与单体机械能 kind=problem alt=none cont=none y=0.79-1.00
\begin{problem}
%FIG p078_g 0.05 0.775 0.225 0.895
\notefig[0.3]{figures/p078_g.png}
静放后，$a$ 向下，$b$ 未离开桌面．
\begin{solution}
\begin{align*}
&W_{Ga}-W_T=\Delta E_{ka}\\
&W_T+W_Q=\Delta E_{kb}\\
&-W_T=\Delta E_{\text{机}a}\\
&W_T=\Delta E_{kb}+Q_f=\Delta E_{\text{机}b}+Q_f\qquad \Delta E_{\text{机}a}\ne\Delta E_{\text{机}b}\\
&W_{Ga}+W_f=\Delta E_{ka}+\Delta E_{kb}\qquad W_{Ga}\to\Delta E_k+Q
\end{align*}
绳拉力对 $ab$ 总功为0
% CHECK: 第2式 W_Q 的下标手写像 Q（按物理应为 f）
\end{solution}
\end{problem}
%%%END
%%%BLOCK id=083-06 ch=06 sec=机械能·能量关系 kind=formula alt=none cont=none y=0.05-0.12
$Q=E_2-E_1$\qquad $E_k=\dfrac12mv^2$

$E_{\text{机}}=E_k+E_p$
%%%END
%%%BLOCK id=083-09 ch=06 sec=重力势能·系统 kind=knowledge alt=none cont=none y=0.20-0.23
$m$与地球系统重力势能$=E_{pm}$
%%%END
%%%BLOCK id=083-12 ch=06 sec=功能关系 kind=formula alt=05,09,12 cont=none y=0.41-0.60
\textbf{功能关系}

% 原稿左侧红笔大括号；“=−mgΔh”“=−kQq/r”写在W_G、W_电式的上方右侧
\[
\left\{\begin{aligned}
&W_{\text{合}}=\Delta E_k\qquad W_G=-\Delta E_{pG}=-mg\Delta h\qquad W_{\text{电}}=-\Delta E_{p\text{电}}=-\frac{kQq}{r}\\[6pt]
&W_{\text{弹}}=-\Delta E_{p\text{弹}}=-\frac12k\Delta x^2\\[6pt]
&W_{\text{引}}=-\Delta E_{p\text{引}}=\frac{GMm}{r} % CHECK: 下标笔迹似“31”，按“引”(万有引力)转录
\\[6pt]
&W_f=-Q_f\\[6pt]
&W_{\text{安}}=-Q_{\text{电}}\\[6pt]
&W_{\text{其他}}=\Delta E_{\text{机}}
\end{aligned}\right.
\]

\[
W_F+W_G+W_{\text{电}}+W_{\text{弹}}+W_f=\Delta E_k
\]
%%%END
%%%BLOCK id=083-13 ch=06 sec=能量图象·有无阻力 kind=method alt=01 cont=none y=0.40-0.58
% 左上一小图：“O”(小球)，下方“↓ ↑”箭头与横线
%FIG p083_b 0.055 0.42 0.16 0.52
\notefig[0.25]{figures/p083_b.png}
%FIG p083_c 0.195 0.405 0.61 0.552
\notefig[0.55]{figures/p083_c.png}
无$f$时

$E_{k1}=mgh$\qquad\qquad 有$f$时\quad 斜率变化
%%%END
%%%BLOCK id=083-14 ch=06 sec=能量图象·斜面上滑下滑 kind=method alt=03 cont=none y=0.58-0.71
%FIG p083_d 0.085 0.589 0.875 0.6805
\notefig[0.85]{figures/p083_d.png}
先\unclear{加}后减速\unclear{来}\qquad $E_k$开口全向上，$E_p$开口全向下。$E_{\text{机}}$-$X$不分段 % CHECK: 首句字迹潦草，“先加(口)后减速来”照录
%%%END
%%%BLOCK id=083-16 ch=06 sec=机车启动 kind=method alt=03 cont=none y=0.85-1.00
\textbf{机车启动}

① 恒$P$启动\qquad $\dfrac{P}{v}-f=ma$\qquad $V_m=\dfrac{P}{f}$
%FIG p083_g 0.125 0.962 0.27 1.0
\notefig[0.25]{figures/p083_g.png}
%FIG p083_h 0.37 0.898 0.525 0.99
\notefig[0.3]{figures/p083_h.png}

② 恒$a$启动
\[
\left\{\begin{aligned}
&F-f=ma\\
&P=Fv
\end{aligned}\right.
\]
\[
P=ma+f\cdot\unclear{at}
\]
% CHECK: 末行被页面下缘截断，仅见“P=ma+f·?t”，原稿似无括号
%FIG p083_i 0.745 0.88 0.935 0.995
\notefig[0.3]{figures/p083_i.png}
%%%END
%%%BLOCK id=084-01 ch=06 sec=机车启动 kind=derivation alt=03 cont=none y=0.03-0.28
\textbf{机车启动方程}

% 原稿左侧有大括号括住下列四式
\[
\left\{\begin{aligned}
&V_m=\frac{P}{f}\\[6pt]
&\frac{P}{at_{\text{加}}}=f+ma\\[6pt]
&\frac12mV_m^2-\frac12m(at_{\text{加}})^2=P(t-t_{\text{加}})-f\Bigl(S_{\text{总}}-\frac12at_{\text{加}}^{2}\Bigr)\\[6pt]
&\frac{P(2t-t_{\text{加}})}{2}=fS_{\text{总}}+\frac12mV_m^2
\end{aligned}\right.
\]
%FIG p084_a 0.608 0.095 0.83 0.25
\notefig[0.3]{figures/p084_a.png}
$F_{\text{恒定}}$\qquad $F\downarrow\to F=f$
%%%END
%%%BLOCK id=084-02 ch=06 sec=机车启动 kind=method alt=03 cont=none y=0.29-0.46
%FIG p084_b 0.11 0.295 0.37 0.45
\notefig[0.35]{figures/p084_b.png}
A$\to$B\quad 恒$a$

B$\to$C\quad 恒$P$

\[
F=\frac{P}{V}-f
\]
%%%END
%%%BLOCK id=084-03 ch=06 sec=机车启动 kind=method alt=03 cont=none y=0.46-0.59
%FIG p084_c 0.085 0.462 0.335 0.585
\notefig[0.3]{figures/p084_c.png}
%FIG p084_d 0.356 0.462 0.575 0.585
\notefig[0.28]{figures/p084_d.png}
\begin{align*}
&\frac{P}{v}-f=ma\\
&a=\frac{P}{m}\cdot\frac1v-\frac{f}{m}
\end{align*}
%%%END
%%%BLOCK id=084-04 ch=06 sec=机车启动 kind=derivation alt=03 cont=none y=0.63-0.90
%FIG p084_e 0.095 0.635 0.30 0.757
\notefig[0.3]{figures/p084_e.png}
\[
f=\frac{F_1V_1}{V_3}\qquad a=\frac{F_1-f}{m}=\frac{F_1-\dfrac{FV_1}{V_2}}{m}
\]
% CHECK: 嵌套分式分母笔迹似V_2(按f=F_1V_1/V_3应为V_3)，分子“F”应为F_1，照原样转录
% 下一式中，两个“=”之间有一弧线箭头，自上一行a的分母m处引下(似表示代入a)
\[
I_{\text{加}}=F_1\frac{V_1}{a}=\ \swarrow\ =\frac{mV_1V_3}{V_3-V_1}
\]
\[
F_1V_1=\frac{F_1}{2}V_2\qquad\therefore V_2=2V_1\qquad\text{反比例函数性质}
\]
%%%END
%%%BLOCK id=084-05 ch=06 sec=机车启动 kind=note alt=01 cont=none y=0.92-1.00
%FIG p084_f 0.08 0.92 0.235 1.0
\notefig[0.25]{figures/p084_f.png}
恒$P$比恒$a$走得远\qquad 故走相同距离\ 恒$P$快
%%%END
%%%BLOCK id=089-06 ch=06 sec=功·摩擦力做功 kind=knowledge alt=none cont=none y=0.70-0.78
\begin{itemize}
\item 相互作用力的功是无关的
\item 斜面摩擦 如走平地
\end{itemize}
% 以上两行左侧有大括号
\begin{keybox}
曲面摩擦\quad 槽$\overset{V}{\text{大}}$拱$\overset{V}{\text{小}}$消耗多
\end{keybox}
% 原稿此行被红笔圈出；“大”“小”二字上方各有一个 V（疑为速度）
%%%END
%%%BLOCK id=089-07 ch=06 sec=功率·平均功率 kind=formula alt=none cont=none y=0.765-0.84
%FIG p089_f 0.705 0.765 0.80 0.81
\notefig[0.25]{figures/p089_f.png}
\[ \bar P_G=\frac Wt\qquad P_{G\,\text{瞬}}=FV\cos\alpha \]
% CHECK: 下标 G 手写形似 G，疑为“功/合”
\noindent 对匀变速\quad $P_{t/2}=\dfrac{P_0+P_t}{2}=\bar P$\qquad $W=\dfrac{(P_0+P)t}{2}=\bar Pt$
%%%END
%%%BLOCK id=090-04 ch=06 sec=动能定理·斜面平抛综合 kind=problem alt=04 cont=none y=0.475-0.62
\begin{problem}
%FIG p090_d 0.06 0.478 0.255 0.585
\notefig[0.3]{figures/p090_d.png}
\begin{solution}
\[ \begin{cases} \dfrac12mV_0^2=mg(H-h)-\mu mg\cos\theta\,\dfrac{L}{\cos\theta}\\[6pt] h=\dfrac{g}{2V_0^2}S^2\end{cases} \]
\[ S^2=2h\left[2(H-h)-2\mu L\right] \]
\[ y=hH-h^2-h\mu L=-h^2+(H-\mu L)h \]
% “−hμL” 写在 “hH−h²” 的下方；“(H−h)”中的 h 笔画加粗
\noindent $h=\dfrac{H-\mu L}{2}$ 时取\unclear{最}大值
\end{solution}
\end{problem}
%%%END
%%%BLOCK id=091-02 ch=06 sec=功率·重力瞬时功率极值 kind=method alt=04,90 cont=none y=0.16-0.29
%FIG p091_b 0.075 0.16 0.215 0.29
\notefig[0.25]{figures/p091_b.png}

$P_{G}\ \max$
\[
P_{G}=mgv\cos\theta=mg\sqrt{2gL\sin\theta}\,\cos\theta
\]
当 $\tan\theta=\dfrac{\sqrt{2}}{2}$ 时取最大
%%%END
%%%BLOCK id=091-07 ch=06 sec=功能关系·多次往返路程 kind=problem alt=90 cont=none y=0.65-0.74
%FIG p091_e 0.065 0.672 0.235 0.735
\notefig[0.25]{figures/p091_e.png}

每次过\unclear{O}动能损失 $5\%$\quad 求总路程
\[
H=H+2H(0.95+0.95^{2}+\cdots+0.95^{n})\xrightarrow[n\to\infty]{}39H
\qquad l_{\text{总}}=\frac{H_{\text{总}}}{\sin\theta}=\frac{39H}{\sin\theta}
\]
% CHECK: 第一个式子左边手写为 H（应为 H总？），原样照录
%%%END
%%%BLOCK id=098-02 ch=06 sec=阻力·需考虑的情形 kind=note alt=05 cont=none y=0.09-0.17
天体\ 大气阻力\quad /\quad 空气阻力\quad /\quad \uline{\unclear{找到…不能…的实例}}

要考虑阻力
% 原稿“阻”字上方有一处很小的批注字，疑为“空气”，未能辨清；上面三项之间以长斜线“/”隔开，最后一项下有横线
%%%END
%%%BLOCK id=098-03 ch=06 sec=弹簧模型·机械能守恒 kind=problem alt=03,08 cont=none y=0.24-0.41
②\ 弹簧振子（简谐运动）\unclear{静放}的从$60^\circ\to120^\circ$ A到最低点，$\mu=0$，\quad $E_{p\text{弹}}\max$\quad $V$从最大$\to0$\quad $a_A$竖直向上
% 原稿此行横贯整页宽度；“竖”字上方有一小处笔迹无法辨认；其下分左右两栏
%FIG p098_b 0.035 0.252 0.195 0.347
\notefig[0.25]{figures/p098_b.png}

% 以下为左栏
$E_{kA\,\max}$\quad $a_A=0$

对A\quad $2T\sin\theta_1=mg$

对B\quad $\therefore N=mg+T\sin\theta_1=\dfrac{3}{2}mg$

当A达$E_{k\max}$前，$a$向下加速

$T\sin\theta<\dfrac{1}{2}mg$\quad $\therefore N_B<\dfrac{3}{2}mg$

% 以下为右栏
$V_B=V_C=0$，整体动能与初始相等。

由机械能守恒定律

$\Delta E_{p\text{弹}}=\Delta E_{pG}$

$V_A=0$时 下降至最低点，$\therefore E_{p\text{弹}}$最大\unclear{的}

$\dfrac{\sqrt{3}-1}{2}\,mgl$
%%%END
%%%BLOCK id=099-03 ch=06 sec=弹性绳·E-h图像 kind=problem alt=03 cont=none y=0.34-0.42
③
%FIG p099_c 0.095 0.338 0.255 0.41
\notefig[0.25]{figures/p099_c.png}

弹性绳\quad $L_0$原长

%FIG p099_d 0.295 0.328 0.45 0.412
\notefig[0.25]{figures/p099_d.png}

\red{\struck{$F-mg=ma$}}
% 原稿红笔所写，被红线划去；首字母 F/T 看不清（疑为 F）
%%%END
%%%BLOCK id=110-06 ch=06 sec=弹性势能 kind=derivation alt=none cont=none y=0.87-1.0
\textbf{用弹性势能的动能定理}
\[
E_{p\text{弹}}=\dfrac{1}{2}k\Delta x^2
\]
%FIG p110_f 0.295 0.915 0.44 0.995
\notefig[0.3]{figures/p110_f.png}
\[
W=\dfrac{k(x_1+x_2)(x_2-x_1)}{2}
\]
%%%END
%%%BLOCK id=117-01 ch=06 sec=变质量功能 kind=method alt=none cont=none y=0.03-0.34
\textbf{变质量功能}

① 重心位移法\quad／\quad ② 图象法\quad／\quad ③ 线密度法

%FIG p117_a 0.07 0.13 0.208 0.245
\notefig[0.25]{figures/p117_a.png}

%FIG p117_b 0.19 0.13 0.436 0.245
\notefig[0.3]{figures/p117_b.png}

③
\[ -a m g\,\frac{a}{2} = -L m_0\,\frac{L}{2}\,g + \frac{1}{2}m_0 L v^2 \]% CHECK: 左边 m 疑漏写下标 0（应为 m_0）

② 图像法.
\[ W=\frac{1}{2}\left(mg+\frac{a}{L}ma\right)(L-a)=\frac{1}{2}mv^2 \quad\therefore\ v=\sqrt{\frac{L^2-a^2}{L}\,g} \]% CHECK: 括号内最后一项原稿像 ma（应为 mg）；L^2 与 a^2 之间的减号不够清晰

（法3）加摩擦
\[ -a m_0 g\,\frac{a}{2} = -m_0 L\cdot\frac{L}{2}\,g + \frac{1}{2}m_0 L v^2 + Q \]

有$\mu$
%FIG p117_c 0.15 0.245 0.40 0.345
\notefig[0.28]{figures/p117_c.png}
% 图中标注：纵轴 f；起点处 $\mu\frac{L-a}{L}mg$，或 $\mu m_0 (L-a) g$；横轴终点 $L-a$

\[ W_f=\frac{1}{2}(L-a)\,\mu\, m_0 (L-a) g \]
%%%END
%%%BLOCK id=117-02 ch=06 sec=功能关系·人做功 kind=note alt=none cont=none y=0.36-0.40
人站起来\underline{做功了}\quad \underline{人原地走}\quad $\underline{W_{\text{内}}=W_f}$
%%%END
%%%BLOCK id=119-03 ch=06 sec=杆绳连接体·机械能与动量 kind=problem alt=04,07 cont=none y=0.29-0.62
\begin{problem}[1.]
%FIG p119_c 0.065 0.30 0.255 0.405
\notefig[0.25]{figures/p119_c.png}
% 图中标注：m、L、m（上方）；A、B；L、L；m、C

① 求AB相碰瞬间速度

② AB速度与C到细杆距离的关系
\begin{solution}
动能定理

① \[ mg\left(L-\frac{\sqrt{3}}{2}L\right)=\frac{1}{2}\cdot 2mv^2 \qquad v=\sqrt{\left(1-\frac{\sqrt{3}}{2}\right)gL} \]

水平方向动量守恒\quad $mv_A+mv_B=0$

%FIG p119_d 0.052 0.445 0.15 0.535
\notefig[0.25]{figures/p119_d.png}
% 图中标注：$\theta$、$L$、$h$

② \[ mg\left(h-\frac{\sqrt{3}}{2}L\right)=\frac{1}{2}\cdot 2mv^2+\frac{1}{2}mv_C^2 \]
\[ v\cos\theta=v_C\sin\theta \]
\[ \tan\theta=\frac{h}{\sqrt{L^2-h^2}}\qquad\therefore\ v=\sqrt{\frac{gh^2\left(2h-\sqrt{3}L\right)}{L^2+h^2}} \]
\end{solution}
\end{problem}
%%%END
%%%BLOCK id=119-05 ch=06 sec=杆连三球·V形模型 kind=problem alt=07,03 cont=none y=0.80-1.00
同
%FIG p119_g 0.535 0.795 0.945 0.86
\notefig[0.4]{figures/p119_g.png}
% 图中为四幅小示意图（竖直两球杆、倒V形、水平三球带向上箭头、斜靠直角装置），“同”字在其左侧

%FIG p119_f 0.005 0.805 0.215 0.962
\notefig[0.25]{figures/p119_f.png}
% 图中标注：A、m（顶点）；L、L；B、C 各 m

bc：$v$\ 0\ 有\ 0\quad $W_{\text{正}}$\ $W_{\text{负}}$

$F_{\text{杆}}=0$时\ $a_A=g$\ 其余时刻小于$g$

$v_{A\text{地}}=\sqrt{2gL}$\quad $mgL=\frac{1}{2}mv^2$\quad $v_B=v_C=0$\ 水平动量抵消

A机械能最小时\quad $v_A=0$\quad $F_{\text{杆}}=0$\quad $v_b=v_c=\dfrac{\sqrt{2gL}}{2}$\quad bc对地压力为$mg$

\unclear{有超\ 失重} % 页面底边被截断，字迹不全
%%%END
%%%BLOCK id=139-02 ch=06 sec=弹簧模型·a-x图像 kind=problem alt=05,03 cont=none y=0.44-0.95
\begin{problem}[剪贴印刷题]
在星球 $M$ 上将一轻弹簧竖直固定在水平桌面上，把物体 $P$ 轻放在弹簧上端，$P$ 由静止向下运动，物体的加速度 $a$ 与弹簧的压缩量 $x$ 间的关系如图中实线所示。在另一星球 $N$ 上用完全相同的弹簧，改用物体 $Q$ 完成同样的过程，其 $a$—$x$ 关系如图中虚线所示。假设两星球均为质量均匀分布的球体。\underline{已知星球 $M$ 的半径是星球 $N$ 的 $3$ 倍，则}% 原稿此处有手画下划线
%FIG p139_c 0.31 0.455 0.57 0.637
\notefig[0.4]{figures/p139_c.png}

A. $M$ 与 $N$ 的密度相等

B. $Q$ 的质量是 $P$ 的 $3$ 倍\quad $M_Q=6M_P$% 手写（大写 M）

C. $Q$ 下落过程中的最大动能是 $P$ 的 $4$ 倍

D. $Q$ 下落过程中弹簧的最大压缩量是 $P$ 的 $4$ 倍\quad 对称性 \unclear{$\checkmark$}% 手写

\textbf{手写批注：}A\,\struck{B}\,C% CHECK: 大字手写 A B C，B 上有一竖线疑为划去（标准答案 AC）

$E_{pQ}=\dfrac{2a_0x_0}{2}$，\quad $E_{pP}=\dfrac{3a_0x_0}{2}$% CHECK: 下标原稿形似 pQ、pP，可能是 kQ、kP；图中纵轴 a 前被手写加了 m；图上另有一条红色手绘曲线

面积是能量

$E=max$

\begin{solution}
\textbf{21. AC}\quad 【命题意图】本题考查万有引力定律、牛顿第二定律和功能关系，意在考查考生的推理能力和分析综合能力。

【解题思路】设 $P$、$Q$ 的质量分别为 $m_P$、$m_Q$；$M$、$N$ 的质量分别为 $M_1$、$M_2$，半径分别为 $R_1$、$R_2$，密度分别为 $\rho_1$、$\rho_2$；$M$、$N$ 表面的重力加速度分别为 $g_1$、$g_2$。在星球 $M$ 上，弹簧压缩量为 $0$ 时有 $m_Pg_1=3m_Pa_0$，所以 $g_1=3a_0=G\dfrac{M_1}{R_1^2}$，密度 $\rho_1=\dfrac{M_1}{\frac{4}{3}\pi R_1^3}=\dfrac{9a_0}{4\pi GR_1}$；在星球 $N$ 上，弹簧压缩量为 $0$ 时有 $m_Qg_2=m_Qa_0$，所以 $g_2=a_0=G\dfrac{M_2}{R_2^2}$，密度 $\rho_2=\dfrac{M_2}{\frac{4}{3}\pi R_2^3}=\dfrac{3a_0}{4\pi GR_2}$；因为 $R_1=3R_2$，所以有 $\rho_1=\rho_2$，选项 A 正确。当物体的加速度为 $0$ 时有 $m_Pg_1=3m_Pa_0=kx_0$，$m_Qg_2=m_Qa_0=2kx_0$，解得 $m_Q=6m_P$，选项 B 错误。根据 $a$—$x$ 图线与坐标轴围成图形的面积和质量的乘积表示合外力做的功可知，$E_{\mathrm{km}P}=\dfrac{3}{2}m_Pa_0x_0$，$E_{\mathrm{km}Q}=m_Qa_0x_0$，所以 $E_{\mathrm{km}Q}=4E_{\mathrm{km}P}$，选项 C 正确。根据运动的对称性可知，$Q$ 下落时弹簧的最大压缩量为 $4x_0$，$P$ 下落时弹簧的最大压缩量为 $2x_0$，选项 D 错误。
\end{solution}
\end{problem}
%%%END
%%%BLOCK id=146-02 ch=06 sec=机械能守恒·斜劈与直杆连接体 kind=problem alt=02,03 cont=none y=0.19-0.50
\begin{problem}[1]
\textbf{1、}①对斜面施加\unclear{力}$F_1$使之平衡，求$F_1$。②用$F$使斜劈缓慢右移，$m$上升$h$，求斜劈位移$S_1$及$F$做的功。% h 后的“,”按文意应为逗号；下文记作 h_1
%FIG p146_b 0.105 0.215 0.26 0.345
\notefig[0.25]{figures/p146_b.png}
\begin{solution}
\[ \begin{cases}
N\cos\theta=mg\\
F_1=N'\sin\theta\\
N'=N
\end{cases} \]
\[ F_1=mg\tan\theta \]
\[ V_{\text{杆}}=V_{\text{斜}}=0 \]
\[ W=mgh_1\qquad S_1=\dfrac{h_1}{\tan\theta} \]
③初始时$B$点高$h_2$，由静止释放$m$、$M$，求$m$着地时$M$速度
\[ mgh_2=\dfrac{1}{2}mV_1^2+\dfrac{1}{2}MV_2^2 \]
\[ V_1=V_2\tan\theta \]
\[ V_2=\sqrt{\dfrac{2mgh_2}{m\tan^2\theta+M}} \]
%FIG p146_c 0.635 0.415 0.74 0.455
\notefig[0.15]{figures/p146_c.png}
\[ \tan\theta=\dfrac{S_{\text{杆}}}{S_{\text{斜面}}} \]
\end{solution}
\end{problem}
%%%END
%%%BLOCK id=146-03 ch=06 sec=功能关系·弹簧与轻绳系统 kind=problem alt=08,03 cont=none y=0.50-0.715
\begin{problem}[2]
\textbf{2、}$\because$简谐运动对称性　$\angle A$从$120^\circ\to60^\circ$过程
%FIG p146_d 0.115 0.514 0.25 0.605
\notefig[0.2]{figures/p146_d.png}
\begin{solution}
$A$刚开始的$a_A=g$向下。

故$A$在最低点时$a_A=g$向上

$\therefore N_{\text{杆}A}=2mg$
\[ 2T\cos30^\circ-mg=ma\qquad T=\dfrac{2\sqrt{3}}{3}mg \]
\[ E_p=mg(h_1-h_2)=\dfrac{\sqrt{3}-1}{2}mgL \]
\[ W_{\text{轻绳}B}=\dfrac{\sqrt{3}-1}{4}mgL\qquad W_{\text{轻绳}A}=\dfrac{\sqrt{3}-1}{4}mgL \]
$P_{\text{绳}B}$先$\uparrow$后$\downarrow$　　$V$为0有0
\end{solution}
\end{problem}
%%%END
%%%BLOCK id=147-01 ch=06 sec=功能关系·斜面弹簧模型 kind=problem alt=none cont=none y=0.02-0.12
% 原稿无题干文字，仅有图和 A、B、C 三行(疑为选择题选项)，B、C 行右侧各有红色勾号
\begin{problem}
%FIG p147_a 0.05 0.02 0.24 0.105
\notefig[0.25]{figures/p147_a.png}
\noindent A：B动能增加量$=$斜面支持力和弹簧弹力做功之和　\red{$W_{\text{合}}=W_{G}+W_{\text{弹}}+W_{N}$}% CHECK: A行按原稿读作“斜面支持力和弹簧弹力做功之和”(未见重力项)，其右侧红笔另批 W合=W_G+W_弹+W_N(下标 G、弹 笔迹较糊)

\noindent B：B机械能增量等于斜面支持力和弹簧弹力做功之和　\red{$\checkmark$}

\noindent C：B和弹组系机增$=$斜面支持力和弹簧弹力做功之和　\red{$\checkmark$}% CHECK: C行左半为缩写，疑为“B和弹簧组成系统机械能增量”；“系”字原稿写成与“系统(孑统)”相同的字形
\end{problem}
%%%END
%%%BLOCK id=147-02 ch=06 sec=动能定理·系统与内力做功 kind=knowledge alt=03,07 cont=none y=0.12-0.37
% 原稿“区别”部分与左侧①②③之间有一条竖线；“系统”原稿写作“孑统”
\textbf{绳杆连接体}

\noindent ①　系统牛二律　$\sum F=\sum(ma)$

\noindent ②　系统　动量定理　$\sum I_{\text{外}}=\Delta\bigl(\sum p_{\text{系}}\bigr)$

\noindent ③　系统　动能定理　$\sum W_{\text{外}}+\sum W_{\text{内}}=\Delta\bigl(E_{\mathrm{k}\text{系}}\bigr)$% CHECK: “=Δ”处被引线笔迹盖住，按上下文读作 =Δ；E_k 的下标读作“系”

\medskip
\noindent \textbf{区别}　①　不考虑内力　\struck{\unclear{②}}% ①下方有一个被涂划掉的带圈符号，辨认不清(疑为②)

\noindent \hspace*{3.2em}②　要考虑内力　不能忽略电场力

\noindent \hspace*{6em}内力$\times$相对距离 $\left\{\begin{tabular}{@{}l@{}}安培力\\摩擦力\\弹簧\end{tabular}\right.$
%%%END
%%%BLOCK id=147-03 ch=06 sec=动能定理·系统与内力做功 kind=note alt=none cont=none y=0.27-0.33
% “C”形符号按“⊂”转录；机守=机械能守恒，能守=能量守恒(原稿缩写)
机守$\subset$能守$\subset$动能定理
%%%END
%%%BLOCK id=147-04 ch=06 sec=动能定理·绳杆连接体 kind=method alt=04 cont=none y=0.35-0.47
\textbf{绳杆\unclear{双}动能定理}% “杆”与“动能定理”之间的字笔迹难辨，疑为“双”

\noindent $\left\{\begin{tabular}{@{}l@{}}1：物1　动能定理\\2：物2　动能定理\\关联速度　$V_1$与$V_2$关系（沿绳/杆速度都相等）\end{tabular}\right.$
%%%END
%%%BLOCK id=164-06 ch=06 sec=机械能守恒·动能-位置图像 kind=method alt=04 cont=none y=0.86-1.00
%FIG p164_g 0.04 0.885 0.19 0.96
\notefig[0.2]{figures/p164_g.png}
\[ mgh=E_K \]
\[ h^2=R^2-(R-x)^2 \]
\[ h=\sqrt{R^2-(R-x)^2}\qquad\text{椭圆} \]
%FIG p164_h 0.42 0.875 0.63 0.985
\notefig[0.3]{figures/p164_h.png}
%%%END
